\ex{讲义例题\quad 弹簧、质量体与阻尼器（第30--33页）}{
质量为 $m$ 的小车沿水平轨道运动，两侧分别通过刚度为 $k$ 的弹簧和黏性阻尼系数为 $f$ 的阻尼器与固定支座相连。外力 $u(t)$ 与位移 $y(t)$ 的正方向均向右。力施加前小车在静止平衡位置，$y$ 从该位置计量。建立以 $u$ 为输入、$y$ 为输出的状态空间模型。
\begin{center}
\begin{tikzpicture}[>=Stealth]
    \draw[gray] (-3.15,.92) |- (3.15,-.45) -- (3.15,.92);
    \foreach \x in {-3.15,3.15} {
        \foreach \h in {0,.21,.42,.63,.84} {\draw[gray] (\x,\h) -- ++(-.15,-.15);}
    }
    \draw[thick,fill=mytheorembg] (-.68,-.14) rectangle (.68,.75);
    \node at (0,.31) {$m$};
    \draw[thick] (-.35,-.29) circle[radius=.15] (.35,-.29) circle[radius=.15];
    \draw[myb,thick] (-3.15,.31) -- (-2.7,.31) --
       (-2.54,.49) -- (-2.32,.13) -- (-2.1,.49) -- (-1.88,.13) --
       (-1.66,.49) -- (-1.44,.13) -- (-1.28,.31) -- (-.68,.31);
    \node[below,myb] at (-2.09,.13) {$k$};
    \draw[myb,thick] (.68,.31) -- (1.73,.31) (1.73,.09) -- (1.73,.53)
       (1.55,.01) -- (2.17,.01) -- (2.17,.61) -- (1.55,.61)
       (2.17,.31) -- (3.15,.31);
    \node[below,myb] at (1.92,0) {$f$};
    \draw[->,thick,myr] (-1.5,.98) -- node[above] {$u(t)$} (-.65,.98);
    \draw[->] (.68,1.16) -- node[above] {$y(t)$} (1.8,1.16);
\end{tikzpicture}
\end{center}
}

\nt{\textbf{知识点}\quad 牛顿第二定律、黏性阻尼、用位移与速度建立二阶系统。}

\sol 沿 $y$ 的正方向列力平衡。外力为 $u$，弹簧的回复力是 $-ky$，阻尼力是 $-f\dot y$，所以
\[
    m\ddot y=u-f\dot y-ky .
\]
取 $x_1=y$、$x_2=\dot y$，得到 $\dot x_1=x_2$、$\dot x_2=-(k/m)x_1-(f/m)x_2+u/m$。于是
\[
    \dot x=\begin{bmatrix}0&1\\-k/m&-f/m\end{bmatrix}x
    +\begin{bmatrix}0\\1/m\end{bmatrix}u,\qquad
    y=\begin{bmatrix}1&0\end{bmatrix}x,\qquad D=0 .
\]
题设的静止平衡初态对应 $x(0)=0$。消去 $x_2$ 可回到原动力学式。若改为测加速度，输出是
$\ddot y=-(k/m)x_1-(f/m)x_2+u/m$，此时 $A,B$ 相同，输出矩阵改为 $C=[-k/m\ \ {-f/m}]$，且 $D=1/m$。

\begin{note}
同一物理系统可测位移、速度或加速度。先由机理定出 $A,B$，然后按题目的测量量写 $C,D$。常见平衡规律有电路的电压/电流定律、机械系统的力/力矩平衡，以及流体和热过程中的物料/能量守恒。
\end{note}
